% Block A Results. Working draft: analysis/block_a/RESULTS_DRAFT.md
% Every number traces to data/block_a/; analysis/block_a/verify_claims.py
% re-derives the checkable ones. [PI] marks a decision recorded in
% analysis/block_a/DISCREPANCY_REPORT.md.

\section*{Results}

\subsection*{The instrument, and its calibration}

Before applying them to predictions, both measurements were run on the deposited
reference structures, where the answer is already known. Across 167 unique
reference entries the two-instrument predicate returned the expected call on 159
and deviated on nine --- five annotated active, four annotated inactive.
\textbf{This is a derivation-set figure, not a held-out one}: the thresholds
were fitted on 80 of these same reference rows, so 159/167 measures internal
consistency rather than generalisation.

The deviations are informative rather than a nuisance. Several entries annotated
active were solved with an agonist bound but no transducer, and these never form
the NPxxY network: the hydrogen-bond arrangement requires G$\alpha$ clamping, so
the predicate declines to call them active. That reading is supported
directly by experiment: the TM5/TM6 movement ``is observed only when both the
agonist and the G protein or mimetic Nb are present; no movement of TM5, TM6 is
observed when only the agonist is present'' \citep[p.~21]{georgiou2025heterogeneity}.
Others reflect construct
engineering inside the measurement window rather than receptor conformation.
Five of the nine deviations are unclassified as to cause; the remainder split
between expected biology, curation error, a measurement artifact, and one entry
that could be either. We report the unclassified fraction rather than resolving
it by assertion.

The practical consequence is that ``active structure'' is not one thing. A panel
built on the assumption that every entry annotated active is geometrically
equivalent would inherit that heterogeneity silently (Fig.~1).

\subsection*{Partner-conditioned activation across four backbones}

Given the receptor sequence alone, predictions are predominantly inactive ---
the collapse onto a single basin, and its state, are established for this
receptor class \citep{heo2022multistate,chib2025gpcrstates}. Given
the receptor sequence together with its cognate G$\alpha$, they are
predominantly active. The direction is consistent on all four backbones across
the 40 Class~A receptors (Fig.~2, Table~2), and it behaves as a switch rather
than a graded shift: 111 of 160 apo cells never fire the predicate on any of 25
seeds, while 108 of 159 cognate cells fire on all 25. The distributions behind
those calls, and the shift in magnitude rather than in call, are in
Fig.~S11.

The median outward TM6 tilt shift on adding the partner is
$+5.04$~\AA{} (Boltz-2, 95\% CI $[3.53, 5.55]$),
$+1.05$~\AA{} (Chai-1, $[0.36, 3.76]$),
$+4.81$~\AA{} (OpenFold3, $[3.95, 5.21]$) and
$+5.31$~\AA{} (Protenix2, $[4.63, 5.74]$).
Pooled row-level predicate-active rates move from 16.0\% to 83.4\% (Boltz-2),
32.7\% to 75.0\% (Chai-1), 24.2\% to 79.9\% (OpenFold3) and 14.5\% to 87.1\%
(Protenix2). All intervals are cluster bootstraps over 26 paralog clusters and
are the authoritative ones; the cluster interval is up to $2.2\times$ wider than
the receptor interval. Chai-1 is the soft predictor throughout --- a smaller
shift with a much wider interval --- which recurs across metrics.

\paragraph{Relation to prior work.}
That a co-folded G protein drives an active-state receptor is not new.
\citet{chiesa2025templatebias} showed on 63 post-cutoff Class~A receptor--G$_s$
pairs that ``the models explicitly modeling G-protein binding'' outperform
template-biased alternatives (p.~6302), and attribute the gain to co-folding
rather than to training-set composition (p.~6305);
\citet{ye2026multistatebias} report that ``large protein partners drive clear
conformational switching between states'' where small molecules do not (p.~2).
Our contribution is not the direction of the effect but its scale and its
controls. Four differences are checkable. First, neither study uses an
AF3-lineage backbone; \citet{chiesa2025templatebias} run AlphaFold2 and
AlphaFold-Multimer only, and propose their benchmark as ``a baseline'' for newer
models (p.~6299). Second, neither has an apo arm: only 16 of 145 structures in
that benchmark represent the receptor bound to G protein alone (p.~6302), and in
\citet{zhang2026generalization} the ligand is present in both arms, so the
reported gain is the marginal effect of adding a transducer to an
already-agonist-bound receptor. Third, both score by RMSD to the deposited
active reference of the same complex, which cannot be evaluated where no active
structure exists; we apply an operationalised predicate instead. Fourth, neither
runs a decoy or scrambled-partner arm, so partner occupancy and partner identity
are never separated.

A contrary result deserves naming. \citet{ku2026promise} find partner-induced
prediction failures substantially more common than ligand-induced ones, and
symmetric rather than one-directional. Three things separate that setting from
ours, and the first is decisive: \textbf{their benchmark excludes the motion we
measure, by construction, and they say so}. Their pair-extraction criterion
``preferentially captures large conformational changes'', with the consequence
that ``localized but biologically important motions, such as GPCR TM6
displacement or transporter pocket rearrangements, may be excluded''
\citep[p.~9]{ku2026promise}. Their partner is also any protein chain forming an
assembly rather than a cognate transducer, and their success criterion is
two-sided --- recovering every known state in a cluster --- which is a strictly
harder target than the directional question asked here.

The apo arm is a computational reference condition, not a physical state: no
ligand, no membrane, no basal-activity equivalent. It should not be read as a
physiological inactive receptor. In the cognate arm the G$\alpha$ is physically
docked rather than merely present: 98.2\% of samples meet the interface-contact
criterion, with a median of 66 receptor--G$\alpha$ contacts.

\subsection*{An orthogonal signature corroborates the direction}

The obvious objection is that two thresholds were chosen and the models happened
to cross them. Predictions were therefore additionally scored on the
P5.50--F6.44 PIF connector --- a different structural rearrangement, in a
different region, never used to set a threshold or to call anything (Fig.~S14).

The strong result is agreement at row level. Of predictions the predicate called
active, 204 of 256 (79.7\%) fall below the reference-inactive median on the
connector; of predictions it called inactive, 205 of 256 (80.1\%) fall above the
reference-active median. The direction matches the references on all four
backbones.

The magnitude is roughly one third of the reference scale (median prediction
delta $-0.56$~\AA{} against a reference delta of $-1.51$~\AA). The interval on
that delta, $[-1.20, +0.03]$, includes zero, as do all four per-backbone
intervals. We therefore report the connector as corroborating in direction and
in row-level agreement, and explicitly not as a signed magnitude effect.

\subsection*{Receptor-specific activation amplitude is not reproduced}

Real receptors do not all activate by the same amount; the allosteric coupling
between the orthosteric site and the cytoplasmic ends of TM5 and TM6 differs
markedly between receptors \citep[p.~8]{georgiou2025heterogeneity}. If a model
had learned
activation rather than an active-looking shape, receptors whose reference
structures differ greatly should receive large predicted shifts and those whose
references differ little should receive small ones.

On the NPxxY axis --- the well-powered one, with a reference separation SD of
5.22~\AA{} --- the regression of predicted shift on reference separation gives
slopes of 0.04 (Boltz-2), 0.37 (Chai-1), 0.18 (OpenFold3) and 0.26 (Protenix2)
across 28 Class~A receptors (Fig.~3, Table~3). Three of four show no evidence of
amplitude reproduction. Protenix2's interval, $[0.07, 0.55]$, excludes zero: a
weak signed positive slope, roughly a quarter of the expected
receptor-to-receptor scaling, and far below unity. It excludes zero under all
three inclusion sets, so this is not an artefact of the restriction choice. No
backbone approaches a slope of one.

The tilt axis cannot answer the question at all. Its reference separation SD is
1.17~\AA{} against NPxxY's 5.22~\AA, and under Class-A restriction two of four
slopes are negative (Boltz-2 $-0.30$, Chai-1 $-0.66$). A negative slope is not
physically interpretable: it would mean a model moves receptors backwards in
proportion to how far they should move. Combined with the restricted range, this
is the signature of a predictor with essentially no dynamic range. We therefore
report no tilt amplitude slope in either direction and describe the axis as
uninformative for amplitude --- a property of the instrument, not a result about
the models.

This is the load-bearing negative result, and it must be read alongside the
fraction in Table~2, a median over 39 receptors. \textbf{The fraction is a mean;
the regression is a covariance.} A high mean shift with near-zero covariance
against reference separation is exactly what this corpus shows, and the two are
not in conflict.

\subsection*{Confidence is informative only at the state-defining residues}

pLDDT averaged over the whole complex is diluted by the G$\alpha$ chain, by loops
and by termini. Restricted to the six state-defining anchor residues (3.50,
3.51, 5.58, 6.30, 6.34, 7.53) it carries information about conformational
correctness --- on two of four backbones (Fig.~4, Table~4).

Whole-complex against anchor-restricted Pearson correlations are $-0.10$ (null)
against $-0.22$ (signed) for Boltz-2; $+0.07$ (null) against $-0.16$ (null) for
Chai-1; $-0.26$ (signed) against $-0.63$ (signed) for OpenFold3; and $+0.33$
(signed positive) against $+0.07$ (null) for Protenix2.

The Protenix2 row is the finding. Read over the whole complex its confidence
correlates positively with distance from the active state --- apparent
overconfidence --- and that reversal disappears entirely at anchor grain.
OpenFold3 moves the other way, strengthening from $-0.26$ to $-0.63$. Reporting
a single pooled pLDDT correlation would have produced two opposite and equally
misleading conclusions.

Confidence scores are not uninformative in general, and the distinction matters.
\citet{zhang2026generalization} report that interface predicted TM-score is ``a
useful decision metric'' for \emph{ligand pose} ($R^2 = 0.49$); their
pLDDT-against-receptor-C$\alpha$ result, $R^2 = 0.046$, is the same null we
report here. What a confidence score predicts well and what it predicts badly
are different questions, and pose accuracy is not conformational state. Others
find no usable relationship between pLDDT and conformation
\citep[p.~5]{bryant2024cfold}, and the same failure appears for peptide
placement \citep{junker2026peptidedesign} and for kinase state
\citep{sun2026kinconfbench}.

This is a finding about method rather than about these four models. On the
evidence here --- four backbones, one receptor class --- we would put it no
higher than a hypothesis worth testing elsewhere: that a confidence score
aggregated over a whole complex can invert relative to the same score read at
the residues that define the property in question. We are not aware of a prior
report of a confidence metric changing sign under re-aggregation.
